\section{Observations and Sample}
\label{sec:obs}

Our sample is the twelve fast radio bursts detected simultaneously by CHIME/FRB
\citep{CHIMEFRB2018} at $0.4$--$0.8$\,GHz and by DSA-110 \citep{Law2024} at
$\sim$1.4\,GHz between 2022 February and 2024 February (MJD $59617$--$60369$).
The defining asset of the sample is that each burst is seen twice across a
$\sim$1\,GHz frequency baseline, which---as developed in
Section~\ref{sec:jointfit}---breaks the degeneracy between the scattering time
and its spectral index that defeats any single-band fit. The bursts span a
dispersion measure range of $262$--$960$\,pc\,cm$^{-3}$; their nicknames, TNS
designations, epochs, widths, and time-of-arrival residuals are listed in
Table~\ref{tab:toa}. We report all twelve, noting per sightline the level at
which each quantity is constrained rather than imposing a single detection
threshold.

\subsection{Dynamic spectra and reduction}
\label{sec:data}

For each co-detection we analyze a pair of total-intensity dynamic spectra---one
per facility---derived from the CHIME/FRB and DSA-110 baseband recordings
(24 spectra in total; provenance and array geometry in
\texttt{DATA\_SOURCES.md} of the companion \texttt{dsa110-FLITS} archive). Each
spectrum is loaded with its frequency axis standardized to ascending order and
its per-channel bandpass normalized, so the fitted amplitudes are free of the
band-edge rolloff. All fits, quality gates, and figures in this paper are
produced by the open \texttt{dsa110-FLITS} pipeline (\texttt{pipeline/}), which
we reference throughout for reproducibility. We do not catalog the fine
baseband morphology of each burst here; the only morphological distinction the
analysis depends on---whether a burst resolves into more than one temporal
component---is treated where it matters, in the multi-component fit of
Section~\ref{sec:multicomp}.

\subsection{Scattering fits}
\label{sec:obs-scatt}

The joint CHIME--DSA fit yields a $1.4$\,GHz pulse-broadening time
$\tau_{1\,\mathrm{GHz}}$ wherever a scattering tail is resolved, measured against
the NE2001 Galactic prediction along the sightline. We fit the joint two-band
model across the co-detected sample with the turbulence spectral index $\beta$ as
the fundamental parameter, from which both the PBF shape and the frequency-scaling
index $\alpha=2\beta/(\beta-2)$ follow (Section~\ref{sec:jointfit}), modeling
temporal multiplicity per band (Sections~\ref{sec:results-alpha},
\ref{sec:disc-alpha}). Sightlines with interior posteriors return a measured
$(\tau_{1\,\mathrm{GHz}},\beta)$ pair; sightlines whose posterior rails at the
$\beta=4$ bound are consistent with the square-law (exponential-PBF) limit and
are quoted only as geometry-conditioned $\alpha=4$ limits, never as
measurements, while a posterior railing at the steep bound or failing the
quality gate is withheld entirely. The final per-sightline result of this
uniform re-fit is reported in Table~\ref{tab:beta} (eleven of twelve
sightlines; one gate failure), and Figure~\ref{fig:budget} overlays every
measured $\tau_{1\,\mathrm{GHz}}$ on the predicted intervening contribution
along its sightline (Section~\ref{sec:budget}).
Figure~\ref{fig:jointmodel_montage} shows the recovered joint model against the
DSA dynamic spectrum for every co-detection. Panels are ordered by burst
nickname; each pair is [data~|~model] on a common color scale, with the fitted
$\beta$ and DSA $\chi^2_\nu$ annotated. FRB~20220310F retains the dedicated
multiplicity demonstration of Figure~\ref{fig:whitney_mult}.

\begin{figure*}[!t]
  \centering
  \includegraphics[width=\textwidth]{jointmodel_montage}
  \caption{Joint-fit model vs.\ data on the DSA-110 band for all twelve
  co-detections. Each row pair shows the on-pulse dynamic spectrum (left) and the
  gain-marginal model reconstruction (right) from the uniform $\beta$-based
  two-band fit (PBF shape and frequency scaling both derived from the fitted
  turbulence index $\beta$; Section~\ref{sec:jointfit}), at the per-burst
  component multiplicities of Table~\ref{tab:beta}. The fitted $\beta$ and
  DSA reduced $\chi^2$ are listed per panel; FRB~20220310F is shown with its
  committed two-component (C2D2) fit, whose multiplicity dependence is
  demonstrated in Figure~\ref{fig:whitney_mult}, and FRB~20240203A's flat
  model panel against strongly structured data visualizes its quality-gate
  failure (Table~\ref{tab:beta}). Rendered from the campaign joint fits by
  \texttt{dsa110-FLITS}
  (\texttt{analysis/scattering-refit-2026-06/plot\_jointmodel\_montage.py}).}
  \label{fig:jointmodel_montage}
\end{figure*}

\subsection{Foreground-galaxy search}
\label{sec:obs-fg}

To establish which sightlines could carry an intervening dispersion or
scattering contribution, we searched each of the twelve fields for foreground
galaxies and clusters, cross-matching candidates against the DESI Legacy Imaging
Survey DR9 with \citet{Zhou2021} photometric redshifts, DESI DR1 spectroscopic
redshifts, the NASA/IPAC Extragalactic Database, and PS1-STRM \citep{Beck2021}.
The search returns 49 candidate systems---34 galaxy halos and 15 clusters---along
nine of the twelve sightlines (three fields yield no candidate). For each
candidate we record the proper impact parameter (12--280\,kpc for the halos;
clusters by $b/R_{500}$) and the best available redshift and its provenance
(spectroscopic or photometric); the full catalog is Table~\ref{tab:foreground}.
Spectroscopic host redshifts are available for nine of the twelve bursts
($z=0.043$--$0.51$), while three (FRB 20230325A, FRB 20240122A, FRB 20230814B) carry only a
placeholder redshift and are excluded from any distance-dependent quantity. We measure no host stellar masses; the only masses entering the
analysis are the halo-mass proxies of the intervening systems used by the
circumgalactic dispersion model (Section~\ref{sec:dm}), whose confidence we flag
per system. The verdict logic that classifies each candidate as foreground,
background, or inconclusive, and the resulting dispersion budget, follow in
Section~\ref{sec:budget}.
